%% The Wind-AE windsoln / seed file format.
%% Author: Kwang-Il Seon
\documentclass[english,a4paper]{my_memo}
\synctex=1
\usepackage{booktabs}
\usepackage{array}
\usepackage{amsmath}
\usepackage{babel}
\emergencystretch=3em

\newcommand{\code}[1]{\texttt{#1}}
\newcommand{\Rp}{R_{\rm pl}}
\newcommand{\Ftot}{F_{\rm tot}}
\newcommand{\hd}[1]{\code{\##1}}      % a #key header token

\title{The Wind-AE \code{windsoln} (seed) file format}
\author{Kwang-Il Seon}
\date{2026-06-13}

\begin{document}
\maketitle

\begin{quote}\small\itshape
A converged Wind-AE wind solution is stored as a single CSV-like file: a
block of comment-prefixed header lines that carry every parameter,
boundary condition, and the stellar spectrum, followed by the solution
profile itself as a table of normalized (code-unit) columns. The same
format is used by (i) the original Wind-AE Python wrapper (Murray-Clay et
al.\ 2009; Broome et al.\ 2025), (ii) the included solution grid under
\code{inputdata/windae\_grid/}, and (iii) EXHALE's Fortran port, whose
\code{load\_seed} reads such a file as the starting point of a
continuation and whose \code{dump\_seed} writes one back. This note is the
authoritative field-by-field reference, taken from the Python reader
(\code{wrapper/wrapper\_utils/windsoln.py}).
\end{quote}

\section{File layout}
A windsoln file has three kinds of lines, in order:
\begin{enumerate}\itemsep2pt
  \item \textbf{Header lines} \code{\#key: v1,v2,\dots} --- one per key,
        carrying scalars, parameters, BCs, flags, and spectrum metadata
        (Table~\ref{tab:hdr}).
  \item \textbf{Spectrum rows} \code{\#\# e,wPhi,$\sigma_1$,$\sigma_2$}
        --- \code{npts} of them, holding the stellar spectrum sampled at
        the Gauss--Legendre quadrature points (\S\ref{sec:spec}).
  \item \textbf{Data rows} --- the solution profile, one grid point per
        line, columns named by \hd{vars} and normalized by \hd{scales}
        (\S\ref{sec:grid}).
\end{enumerate}
Any other line beginning with \code{\#} is a free comment and is skipped.
A minimal excerpt:
\begin{verbatim}
#nspecies: 2
#vars: r,rho,v,T,Ys_HI,Ys_HeI,Ncol_HI,Ncol_HeI,q,z
#scales: 8.34e9,1.0e-15,9.08e5,1.0e4,1,1,1.0e17,1.0e17,1,8.34e9
#plnt_prms: 2.348e30,8.34e9,1.870e33,4.937e11,2.42e4,3.828e33
#phys_prms: 0.800,0.200,HI,HeI,1.673e-24,6.646e-24,2.30
#bcs: 1.030,7.755,2.466e4,0.198,1.0,1.0,1.08e-3,2.05e-4,1.85e-5,1.50e-5
#tech: 1.0,1.0e2,2.5,0.99
#flags: 1,1.00000,0.00000,1
#add_prms: 0
#spec_prms: 59,2,2009-01-01,scaled-solar,full,5.99e-8,9.12e-6, ... ,2.18e-11,3.94e-11
## 2.150e-9,1.771e6,4.258e-24,1.267e-22
## ... (59 rows) ...
1.030,2.466e4,2.232e-6,0.198,1.0,1.0,7.34e3,4.70e2,0.0,2.76
... (one row per grid point) ...
\end{verbatim}

\section{Header reference}
Table~\ref{tab:hdr} lists every header key. $N$ is the number of species
(\hd{nspecies}; $N=2$ for the included H/He data). Planet parameters are in
cgs; boundary conditions and the data grid are in the code units of
Table~\ref{tab:grid}.

\begin{table}[t]\centering\small
\caption{Header lines of a windsoln file. Field order is exact;
``$\times N$'' marks a block repeated for each species.}
\label{tab:hdr}
\begin{tabular}{@{}l l p{8.4cm}@{}}
\toprule
Key & Fields & Meaning (and units) \\
\midrule
\hd{nspecies}  & $N$ & number of species ($N=2$: H\,\textsc{i}, He\,\textsc{i}). \\
\hd{vars}      & names & data-grid column labels (Table~\ref{tab:grid}). \\
\hd{scales}    & $N_{\rm col}$ & cgs value of one code unit for each data column. \\
\hd{plnt\_prms}& 6 & $M_{\rm p}$[g], $\Rp$[cm], $M_\star$[g], $a$[cm],
                     $\boldsymbol{\Ftot}$\,[erg\,cm$^{-2}$s$^{-1}$],
                     $L_\star$[erg\,s$^{-1}$]. \emph{Field~5 is the total
                     ionizing flux, not $T_{\rm eq}$.} \\
\hd{phys\_prms}& $3N{+}1$ & mass fractions $X_s$ ($\times N$); species names
                     ($\times N$); atomic masses [g] ($\times N$);
                     then \code{molec\_adjust} (mean-molecular-weight boost
                     of the molecular base layer, $\approx2.3$). \\
\hd{bcs}       & $2N{+}6$ & $R_{\rm min},R_{\rm max}$ [$\Rp$];
                     $\rho_{\rm min}$ [$\rho_0$]; $T_{\rm min}$ [$T_0$];
                     base neutral fractions $Y_{s,\rm min}$ ($\times N$);
                     sonic-point columns $N_{{\rm col},s}$ [$N_{\rm col,0}$]
                     ($\times N$); \code{erf\_drop} (2: the erf-velocity
                     molecular$\to$atomic transition). \\
\hd{tech}      & 4 & \code{breezeparam}, \code{rapidity}, \code{erfn},
                     \code{mach\_limit} (relaxation/solver controls). \\
\hd{flags}     & 4 & \code{lyacool} (0/1); \textbf{\code{tidalforce}}
                     (0 = spherical, 1 = tidal/Roche); \textbf{\code{bolo\_heat\_cool}}
                     (0 = atomic base, 1 = molecular bolometric layer);
                     \code{integrate\_outward} (0/1). \\
\hd{add\_prms} & $1{+}2n$ & $n$, then $n$ $(\text{name},\text{value})$ pairs
                     of optional extra parameters ($n=0$: none). \\
\hd{spec\_prms}& meta & \code{npts}, $N$, date, source file
                     (\code{scaled-solar}), kind (\code{full}),
                     window[2], resolved[2], normalized[2] (wavelength
                     bounds [cm]), ionization potentials $\chi_s$ [erg]
                     ($\times N$). See \S\ref{sec:spec}. \\
\bottomrule
\end{tabular}
\end{table}

\section{The embedded stellar spectrum}
\label{sec:spec}
The stellar XUV spectrum that the solution was computed with is stored in
the file itself, in two parts:
\begin{itemize}\itemsep2pt
  \item \hd{spec\_prms} --- metadata: the number of quadrature points
        \code{npts}, the source label (here \code{scaled-solar}, a
        solar-like SED), the spectral kind (\code{full}, i.e.\ resolved
        rather than monochromatic), the integration window in wavelength
        ($\approx[6,912]\,$\AA, i.e.\ X-ray to the H Lyman limit), and the
        the ionization potentials $\chi_s$ for each species (e.g.\
        $2.18\times10^{-11}\,$erg $=13.6$\,eV for H,
        $3.94\times10^{-11}\,$erg $=24.6$\,eV for He).
  \item the \code{\#\#} rows --- the spectrum proper, one row per
        quadrature point: photon energy $hc/\lambda_i$ [erg]; the
        \emph{normalized} weight $w_i\Phi_{\lambda_i}/\Ftot$; and the
        photoionization cross sections $\sigma_{\lambda_i,s}$ [cm$^2$] for
        each species.
\end{itemize}
The crucial point is the separation of \emph{shape} and \emph{amplitude}:
the \code{\#\#} weights are normalized by $\Ftot$, so they encode only the
spectral \emph{shape} (which photon energies carry what fraction of the
ionizing flux), while the absolute level is set by the single scalar
$\Ftot$ in \hd{plnt\_prms}. This is why the included grid can keep one fixed
\code{scaled-solar} shape across all $\sim$1000 solutions and vary only
$\Ftot$ (and $M_{\rm p},\Rp$) from planet to planet. Verified: \hd{spec\_prms}
is byte-identical across the entire grid.

\medskip\noindent\textbf{Caveat (spectrum dependence).} The wind structure
depends on the SED shape, not only on $\Ftot$. The included data assume a
\code{scaled-solar} spectrum; an EXHALE run that instead assumes, say, a
power law produces different heating/ionization profiles. A Wind-AE
warm-start is therefore only self-consistent with an EXHALE run that uses
the \emph{same} spectrum. Conversely, a power-law EXHALE model failing to
converge does \emph{not} imply that no steady wind exists for that planet
--- only that that particular spectrum/numerics combination did not settle.

\section{The solution grid}
\label{sec:grid}
Each data row is one grid point; the columns are named by \hd{vars} and are
stored in dimensionless code units. The physical (cgs) value is recovered
as $\text{(physical)} = \text{(code)}\times\text{(scale)}$, with the scale
taken from \hd{scales} (Table~\ref{tab:grid}). The independent variable is
the relaxation coordinate $q\in[0,1]$ mapping the base to the sonic point;
the physical radius is $r = R_{\rm min} + q\,z$, where $z$ (constant along
the solution) is the base$\to$sonic span in $\Rp$.

\begin{table}[t]\centering\small
\caption{Data-grid columns (\hd{vars}) and their code$\to$cgs scales
(\hd{scales}). $\rho_0{=}10^{-15}$, $c_{s,0}{=}\sqrt{k_BT_0/m_{\rm H}}
{=}9.08\times10^{5}$, $T_0{=}10^{4}$, $N_{\rm col,0}{=}10^{17}$.}
\label{tab:grid}
\begin{tabular}{@{}l l l p{6.0cm}@{}}
\toprule
Column & Symbol & Scale & Meaning \\
\midrule
\code{r}        & $r$            & $\Rp$              & radius \\
\code{rho}      & $\rho$         & $\rho_0$           & mass density [g\,cm$^{-3}$] \\
\code{v}        & $v$            & $c_{s,0}$          & radial velocity [cm\,s$^{-1}$] \\
\code{T}        & $T$            & $T_0$              & temperature [K] \\
\code{Ys\_HI}   & $Y_{\rm HI}$   & 1                  & H neutral fraction \\
\code{Ys\_HeI}  & $Y_{\rm HeI}$  & 1                  & He neutral fraction \\
\code{Ncol\_HI} & $N_{\rm HI}$   & $N_{\rm col,0}$    & H column to the star [cm$^{-2}$] \\
\code{Ncol\_HeI}& $N_{\rm HeI}$  & $N_{\rm col,0}$    & He column to the star [cm$^{-2}$] \\
\code{q}        & $q$            & 1                  & relaxation coordinate $\in[0,1]$ \\
\code{z}        & $z$            & $\Rp$              & base$\to$sonic span ($r=R_{\rm min}+qz$) \\
\bottomrule
\end{tabular}
\end{table}

\section{Use in the EXHALE port}
\code{src/modules/wind\_ae/wae\_continuation.f90} reads and writes this
format:
\begin{itemize}\itemsep2pt
  \item \code{load\_seed} parses \hd{plnt\_prms}, \hd{phys\_prms},
        \hd{bcs}, \hd{tech}, \hd{flags}, and the first $m$ data rows (into
        the relaxation arrays). It \emph{ignores} \hd{scales},
        \hd{add\_prms}, \hd{spec\_prms}, and the \code{\#\#} rows: the
        spectrum is loaded separately from
        \code{inputdata/windae\_spectrum.inp}, which is the same
        \code{scaled-solar} data in the equivalent
        ``\code{\# KEY: \dots}'' header form.
  \item \code{dump\_seed} is the exact inverse: it reconstructs the full
        header (including \hd{scales}, \hd{spec\_prms}, and the
        \code{\#\#} block from the in-memory spectrum arrays) and writes
        the $m$ relaxation rows, so a converged solution can be banked as a
        reusable seed.
  \item \textbf{Seed library and picker.} \code{inputdata/windae\_grid/}
        holds the Broome et al.\ (2025) grid ($\sim$1000 solutions spanning
        $M_{\rm p}\!\in\![3,528]\,M_\oplus$ and
        $\Rp\!\in\![1.9,22.6]\,R_\oplus$, all with \code{tidalforce}=1 and
        the fixed \code{scaled-solar} spectrum), plus the eight
        \code{starting\_solutions}. With \code{Wind-AE seed: grid} in
        \code{input.inp}, \code{pick\_nearest\_seed} reads the grid manifest
        (one row of \code{path}, $M_{\rm p}$, $\Rp$, $\Ftot$ per file;
        rebuilt by \code{make\_windae\_manifest.sh}) and selects the
        solution closest in $(\log M_{\rm p},\log\Rp)$, so the continuation
        starts from a near neighbor rather than the single distant fiducial
        seed. \code{Wind-AE seed out:} dumps the converged solution back
        into this format.
\end{itemize}

\end{document}
